\chapter{TRABAJO FUTURO} \label{TrabajoFuturo}
En este trabajo principalmente se ha realizado una investigación orientada a definir una aplicación de apoyo en la búsqueda y rescate de personas desaparecidas, así como el desarrollo de un prototipo. Esto ha permitido alcanzar un nivel de madurez tecnológica del proyecto de alrededor de un TRL 4, por lo que todavía existe recorrido para conseguir una solución realmente fiable y operativa.

De cara a un futuro próximo, se plantean dos líneas principales de actuación: realizar una primera validación del prototipo desarrollado, dándo a conocer el proyecto y modificando lo necesario para que se adapte lo máximo posible a las necesidades reales; y, una recopilación de datos que permita investigar en profundidad el modelo de predicción propuesto.

Los primeros pasos consistirían en buscar la consolidación del prototipo actual añadiendo las funcionalidades pendientes como el inicio de sesión seguro, los mapas, la geolocalización de participantes en una búsqueda, las notificaciones push o la posibilidad de crear puntos de control en la búsqueda, todo ello implicará crear nuevas tablas o modificar las tablas existentes de la base de datos AlertRes y adaptar el frontend y el backend de la aplicación.

Una vez obtenido un prototipo más robusto y funcional se deberán realizar pruebas de uso más completas con distintos grupos para evaluar la usabilidad, rendimiento, recoger feedback y corregir errores. 

Paralelamente, se deberá recoger u obtener acceso a  conjuntos de datos reales para continuar y profundizar con el estudio del sistema descrito de predicción de la localización de personas desaparecidas en base a las características del caso, de la propia persona y datos históricos. 

En cuanto a las decisiones tecnológicas, en el desarrollo del proyecto se ha optado por utilizar React Native, una tecnología multiplataforma que permite un desarrollo más ágil y sencillo, facilitando la entrega de un prototipo funcional en menos tiempo. Esta elección responde también a la recomendación de los tutores, dado que para un Trabajo Fin de Máster resulta suficiente y adecuado priorizar la rapidez y la simplicidad del proceso. No obstante, React Native presenta limitaciones en rendimiento, acceso a funcionalidades avanzadas del dispositivo y estabilidad en procesos críticos, aspectos muy relevantes en una aplicación destinada a búsqueda y rescate, donde la fiabilidad debe ser máxima. Por ello, se considera que, en un futuro, la evolución natural del proyecto debería orientarse hacia un desarrollo nativo en Android e iOS, garantizando así un mayor control, precisión y robustez en escenarios reales.

A medida que el prototipo alcance un mayor grado de desarrollo, será necesario redefinir el entorno tecnológico, incluyendo las plataformas soportadas, la compatibilidad con APIs externas y la integración con otros sistemas de datos. Asimismo, con el objetivo de garantizar el respaldo institucional y facilitar el acceso a datos reales que impulsen el desarrollo del modelo de predicción, resultará conveniente establecer acuerdos con organismos competentes como el CNDES.

Finalmente, una vez alcanzada una madurez técnica cercana a TRL 9, se podría plantear la salida al mercado de la aplicación. En esta fase, también se contemplaría la incorporación de nuevas funcionalidades orientadas a mejorar la eficacia operativa, como algoritmos capaces de sugerir estrategias óptimas de búsqueda en función de las características del caso, la integración de dispositivos de monitorización de participantes y recursos, o incluso el diseño de sensores específicos para perros de búsqueda cuya información quede registrada y asociada en la aplicación.

Actualmente, el proyecto se enmarca dentro del ámbito de la I+D+i, con potencial de evolucionar hacia una startup. En este sentido, cuenta con oportunidades de financiación a través de diferentes programas de ayudas y subvenciones a nivel nacional y regional (CDTI, ICECYL en Castilla y León o TRADE en Andalucía, entre otros)\cite{icecyl_ayudas, junta_andalucia_trade_idi, cdti_apoyo_idi}, que podrían suponer un apoyo clave para su desarrollo. Asimismo, también se podrían explorar vías de financiación europeas, como el programa Horizonte Europa, que incluye iniciativas como el Clúster 3 de seguridad civil para la sociedad. Este tipo de iniciativas cuenta con algunas convocatorias, como la actualmente abierta (HORIZON-CL3-2026-01-FCT-03)\cite{horizon2026_fct03}, que están orientadas a proyectos relacionados con personas desaparecidas y buscan alcanzar resultados en niveles de madurez tecnológica intermedios (TRL 6-8), lo que encajaría con la evolución prevista del proyecto.